\section{Discussion}\label{sec:discussion}

A fit at a solved example has always been cheap, and it has always learned a check. We have shown when that check is
the algorithm: when the answer is completed with the labels that let a local rule verify it, and when the condition that
orders the answer, invisible at the finished answer, is kept in the rule. Then the fit of a few milliseconds and the
procedure that runs at any size are the same object. Because that object is a small table with no size in it, its
correctness can be proved rather than measured. The cold layer of our earlier paper was built for this. Its reads,
whether a witness exists and how two counts compare, are the conditions that survive a change of size, and its zero
temperature keeps a tie count from leaking into a distance, as it does in CRAM's warm hidden layer.

\textbf{Limitations.} Fifteen CLRS-30 tasks read a primitive or a piece of structure we supply, from a sum to the open
nodes of a depth-first search, and two more a change of units or an index shift; they show what the check does with the
right quantity, not that it finds it. On grids, the coordinate vocabulary
is a prior that CRAM does not use, and the puzzle search is a classical depth-first search.

Two directions are open. One is to search for the completion itself: closures and counts form a small grammar, and a
completion that makes the training answers locally checkable is a target a search can recognise. The other is to put
the check where our earlier layer stands, beside a language model, where a fit of milliseconds could replace hours of
training through loops.
